\documentclass[11pt]{article}
\usepackage[margin=1in]{geometry}
\usepackage{amsmath,amssymb,amsthm,hyperref}
\newtheorem{theorem}{Theorem}
\newcommand{\fr}[1]{\left\{#1\right\}}
\title{A finite certificate against the Beatty logarithm bound by one}
\author{Principia Math In Progress}
\date{29 September 2026}
\begin{document}
\maketitle
\noindent\textbf{Task 1 (MathDB 330198).}
For irrational $a,b>1$, put
\[
 E(a,b)=\log(a/b)-\sum_{n=1}^{\infty}
 \frac{a\fr{(n+1)/a}-b\fr{(n+1)/b}}{n(n+1)}.
\]
Polanco's Conjecture 2 asserts $|E(a,b)|<1$ \cite{polanco}.
The sum converges absolutely, since its numerator has absolute value
less than $\max(a,b)$. The conjecture is false.

\begin{theorem}
For every irrational $0<\varepsilon<1/50$, the parameters
$a=2-\varepsilon$ and $b=1+\varepsilon$ satisfy
\[
 E(a,b)>\frac{95797}{92400}>1.
\]
In particular, $\varepsilon=\sqrt{2}/100$ gives an explicit counterexample.
\end{theorem}
\begin{proof}
Both parameters are irrational and exceed one. For $1\le n\le10$,
\[
 \left\lfloor\frac{n+1}{b}\right\rfloor=n,
 \qquad
 \left\lfloor\frac{n+1}{a}\right\rfloor
   =k_n:=\left\lfloor\frac{n+1}{2}\right\rfloor.
\]
Indeed, $n(1+\varepsilon)<n+1<(n+1)(1+\varepsilon)$.
For the second equality, $k_n(2-\varepsilon)<n+1$ and
$(k_n+1)(2-\varepsilon)>n+1$: the latter follows from
$(k_n+1)\varepsilon<6/50<1$, treating odd and even $n$.
Consequently the numerator in the $n$th summand is
\[
 n-2k_n+\varepsilon(n+k_n)=
 \begin{cases}-1+\varepsilon(n+k_n),&n\text{ odd},\\
 \varepsilon(n+k_n),&n\text{ even}.
 \end{cases}
\]
Writing $S_{10}$ for the first ten terms of the series, exact summation gives
\[
 S_{10}=-C+\varepsilon D,\qquad
 C=\sum_{j=1}^{5}\frac{1}{(2j-1)2j}=\frac{1627}{2520},\qquad
 D=\sum_{n=1}^{10}\frac{n+k_n}{n(n+1)}=\frac{18587}{5544}.
\]
For the remaining tail $T$, each numerator is less than $a$, so
\[
 T< a\sum_{n=11}^{\infty}\frac{1}{n(n+1)}
 =\frac{a}{11}<\frac{2}{11}.
\]
Moreover, $a/b>33/17$. For $x>1$ the elementary inequality
$\log x>2(x-1)/(x+1)$ follows by differentiation: the derivative
of their difference is $(x-1)^2/[x(x+1)^2]>0$, and the difference
vanishes at one. Therefore $\log(a/b)>16/25$. Combining the bounds,
\[
 E(a,b)=\log(a/b)+C-\varepsilon D-T
 >\frac{16}{25}+\frac{1627}{2520}
   -\frac{18587}{277200}-\frac{2}{11}
 =\frac{95797}{92400}=1+\frac{3397}{92400}.
\]
Finally, $1<\sqrt2<2$ makes the displayed explicit choice admissible.
\end{proof}

\noindent\textbf{Scope and verification.}
The target appears on page 4 of arXiv:1503.08512v1 and remains
Conjecture 2 on journal page 104 of the final article. Both formulations
allow arbitrary irrational $a,b>1$; complementarity is not required.
The argument directly treats the defining series and bounds its infinite
tail. It does not depend on the source's other identities or on numerical
approximation. The separate bound $\pi^2/6$ (Conjecture 1, MathDB 330197)
is not resolved here. The accompanying verifier checks the finite
arithmetic and floor bounds with exact rational arithmetic.

\medskip
\noindent\textbf{Proposed follow-up tasks (not claimed solved).}
\begin{enumerate}
\setcounter{enumi}{1}
\item Determine the optimal universal constant $K$ with $|E(a,b)|<K$
for all irrational $a,b>1$, including whether the optimal bound is attained.
\item For $F(a)=\log a-a\sum_{n\ge1}\fr{(n+1)/a}/[n(n+1)]$,
derive its two one-sided limits and jump at every rational $a>1$.
\item Determine a quantitative modulus of continuity of $F$ at irrational
points in terms of their rational approximation properties.
\item Determine whether $|E(a,b)|<1$ holds under the additional
complementarity condition $a^{-1}+b^{-1}=1$.
\end{enumerate}

\begin{thebibliography}{1}
\bibitem{polanco} G. Polanco, \emph{The logarithm of irrational numbers
and Beatty sequences}, Acta Arithmetica \textbf{179} (2017), 101--123,
\href{https://doi.org/10.4064/aa6695-1-2017}{doi:10.4064/aa6695-1-2017}.
Preprint: \href{https://arxiv.org/abs/1503.08512}{arXiv:1503.08512v1} (2015).
\end{thebibliography}
\end{document}
